\section{Por qué la verificación es el cuello de botella de la automejora}

La sesión anterior mostró que un muestreo abundante puede lograr que el conjunto de candidatos incluya más respuestas correctas; pero un sistema real debe elegir un solo resultado entre los candidatos. Robust verification estudia precisamente: \textbf{cómo determinar si la respuesta final es correcta, en qué paso comienza el error, y cómo usar la señal de verificación para la selección durante la inferencia y para el entrenamiento del modelo.}

Esta sesión aborda, en orden cronológico, cuatro generaciones de métodos:

\begin{enumerate}
    \item \textbf{Outcome verifier}: solo determina si la respuesta completa es correcta o incorrecta;
    \item \textbf{Process reward model}: anota el proceso de razonamiento paso a paso;
    \item \textbf{process supervision automática}: sustituye la costosa anotación humana por pasos mediante la tasa de éxito de rollouts;
    \item \textbf{ensamble de verificadores débiles}: combina múltiples verifiers imperfectos para estimar una señal de corrección más confiable.
\end{enumerate}

\begin{importantbox}{Los verificadores tienen dos usos completamente distintos}
En la \textbf{etapa de inferencia}, el verifier se usa para rerank, poda y detención; en la \textbf{etapa de entrenamiento}, el verifier funciona como reward o como filtro de datos que modifica los parámetros del modelo. El error del primero afecta a una sola respuesta, mientras que el sesgo sistemático del segundo será aprendido y amplificado por el modelo, por lo que la robustez exigida a un verifier usado en entrenamiento es mayor.
\end{importantbox}

\subsection{Propiedades de un verifier ideal}

Dado un problema $x$ y una solución candidata $y$, el verifier produce $v(x,y)\in[0,1]$. En el caso ideal:

$$
v(x,y)\approx P(y\text{ correct}\mid x,y).
$$

Además de la precisión, también importan la calibration, la generalización a muestras fuera de distribución, la robustez frente al estilo de la respuesta, el costo computacional y la capacidad de resistir el reward hacking.

\subsection{Resumen de la sección}

La verificación no es un paso accesorio de puntuación, sino la interfaz común entre inference scaling y reinforcement learning. Determina qué trayectorias se seleccionan, qué datos se conservan y qué comportamientos reciben gradiente.

